An important limitation of finite-state language models is that they can definitionally only take into account a finite set of contexts. 
However, human language is very diverse and inventive.
It also contains arbitrarily deep recursive structures which cannot be captured by a finite set of possible histories.
Now we are going a rung higher on the ladder of the Chomsky hierarchy in our exploration of formal language models. 
We are going to consider \defn{context-free languages}, a more general class of languages than the regular languages that are described by finite-state automata.
Luckily, we will see that a lot of the formal machinery we introduce in this section closely follows analogs from the finite-state section and we invite the reader to pay close attention to the parallels.
Similarly to how we weighted a string in a regular language by summing over the weights of the paths with labeled with that string, we will weight strings in context-free languages by summing over analogous structures. 
We will then again be interested in how to compute such weights and when language models modeling context-free languages are tight.

As we will see, FSAs cannot recognize context-free languages.
We will therefore have to resort to more general automata, known as \defn{pushdown automata}, which we define below. 
Before giving a formal treatment of pushdown automata, however, we will discuss an arguably more natural formalism for generating the context-free languages---\defn{context-free grammars}.\footnote{You might wonder what non context-free grammars are. 
A more general class is that of context-sensitive grammars, in which a production rule may be surrounded by a left and right context. 
They are still however a set of restricted cases of general grammars, which are grammars that can emulate Turing machines.}

\subsection{Context-free Grammars}
We start by giving a formal definition of context-free grammars and introducing some notation.

\begin{definition}[Context-free Grammar]
A \defn{context-free grammar} (CFG) is a 4-tuple $\grammar = \cfgtuple$ where $\nonterm$ is a non-empty set of non-terminal symbols, $\alphabet$ is an alphabet of terminal symbols with $\nonterm \cap \alphabet = \emptyset$, $\start\in\nonterm$ is the designated start non-terminal symbol\footnote{As is the case for initial states in FSAs, multiple start symbols could be possible. However we consider only one for the sake of simplicity.} and $\rules$ is the set of production rules, where each rule $\arule \in \rules$ is of the form $\production{X}{\valpha}$ with $X\in\nonterm$ and $\valpha\in\kleene{(\nonterm\cup \alphabet)}$.
\end{definition}


CFGs allow us to generate strings $\str\in\kleene{\alphabet}$ by \emph{applying} production rules on its non-terminals. 
We apply a production rule $\production{X}{\valpha}$ to $X\in\nonterm$ in a rule $\arule$ by taking $X$ on the right-hand side of $\arule$ and replacing it with $\valpha$.\footnote{We say that $\NT{X}$ is on the right-hand side of a rule $\arule$ if $\arule$ takes the form $\arule=(\production{\NTY}{\valpha \, \NT{X} \vgamma})$, where $\valpha,\vgamma\in\kleene{(\nonterm\cup \alphabet)}$. 
We will sometimes refer to $\NT{X}$ as the \emph{head} of the production rule $\production{\NT{X}}{\valpha}$, and the right-hand side $\valpha$ as the \emph{body} of the production rule.} See a formal definition below.

\begin{definition}[Rule Application]
    A production rule $\production{\NTY}{\vbeta}, \vbeta\in\kleene{(\nonterm\cup\alphabet)},$ is \defn{applicable} to $\NTY$ in a rule $\arule$, if $\arule$ takes the form 
    \begin{equation*}
        \production{\NTX}{\valpha \, \NTY \, \vgamma}, \quad \valpha,\vgamma\in\kleene{(\nonterm\cup\alphabet)}.
    \end{equation*}
    The \defn{result of applying} $\production{\NTY}{\vbeta}$ to $\valpha \, \NTY \, \vgamma$ is $\valpha \, \vbeta \, \vgamma$.
\end{definition}

Starting with $\start$, we apply $\production{\start}{\valpha}$ to $\start$ for some $(\production{\start}{\valpha})\in\rules$, then take a non-terminal in $\valpha$\footnote{We will write $X\in\valpha$, which formally means a substring of $\valpha$ with length $1$. Unless otherwise stated, $X$ can be either a non-terminal or a terminal.} and apply a new production rule. To generate a string we follow this procedure until all non-terminal symbols have been transformed into terminal symbols. The resulting string, i.e. its \defn{yield}, will be the string taken by concatenating all terminal symbols read from left to right. 
We define a derivation more formally below.

\begin{definition}[Derivation]
    A \defn{derivation} in a grammar $\grammar$ is a sequence $\valpha_1,\dots,\valpha_m$, where $\valpha_1\in\nonterm$, $\valpha_2,\dots,\valpha_{m-1}\in\kleene{(\nonterm\cup \alphabet)}$ and $\valpha_m\in\kleene{\alphabet}$, in which each $\valpha_{i+1}$ is formed by applying a production rule in $\rules$ to $\valpha_{i}$. 
    
    We say that a derivation is a \defn{left-most derivation} if for each $\valpha_i$, a production rule is applied to its left-most non-terminal to form $\valpha_{i+1}$.  
    We say that a derivation is a \defn{right-most derivation} if for each $\valpha_i$, a production rule is applied to its right-most non-terminal to form $\valpha_{i+1}$.
\end{definition}

We say that $\valpha\in\kleene{(\nonterm\cup \alphabet)}$ is derived from $\NTX\in\nonterm$ if we can apply a finite sequence of production rules to generate $\valpha$ starting from $\NTX$.
We will denote this as $\derives{\NTX}{\valpha}$.
See the following formal definition.
\begin{definition}[Derives]
We have that a substring $\valpha$ of $\valpha'\in\kleene{(\nonterm\cup \alphabet)}$ \defn{derives from} $\NTX\in\nonterm$, denoted $\derives{\NTX}{\valpha}$, if either $\production{\NTX}{\valpha'}$, or $\production{\NTX}{\NTY}$ and $\derives{\NTY}{\valpha'}$.
\end{definition}
The (context-free) language of a CFG $\grammar$ is defined as all the strings $\str\in\kleene{\alphabet}$ that can be derived from the start symbol $\start$ of $\grammar$, or alternatively, the set of all yields possible from derivations in $\grammar$ that start with $\start$. We will denote the language generated by $\grammar$ as $\lang(\grammar)$. 

\begin{definition}[Language of a Grammar]
    The \defn{language} of a context-free grammar $\grammar$ is 
\begin{equation}
    \lang(\grammar)=\{\str\in\kleene{\alphabet}\mid \derives{\start}{\str}\}
\end{equation}
\end{definition}

Note that the same set of production rules applied in a different order may generate different strings. Hence, we cannot represent derivations simply as multisets of production rules. Instead, we represent derivations with \defn{derivation trees} 
(also known as parse trees). 
In a derivation tree $\tree$, the start symbol $\start$ is the root node, each internal node is labeled by a non-terminal symbol and each leaf node is labeled by a terminal symbol. Each edge in the tree corresponds to a production rule, where the parent node is the left-hand side non-terminal symbol and the children are the symbols on the right-hand side. In cases where it is irrelevant to consider the order of the production rules in a derivation tree, we will write $(\production{\NTX}{\valpha})\in\tree$ to refer to specific production rules in the tree -- viewing trees as multisets (or bags) over the production rules they include. We will denote the yield of a tree $\tree$ by $\treeyield{\tree}$.

\begin{example}[Nominal Phrases]\label{ex:cfg}

CFGs are often used to model natural languages. Terminals would then correspond to words in the natural language, strings would be text sequences and non-terminals would be abstractions over words. As an example, consider a grammar $\grammar$ that can generate a couple of nominal phrases. We let $\nonterm=\{\NT{Adj}, \NT{Det}, \NT{N}, \NT{Nominal}, \NT{NP}\}$, $\alphabet=\{\text{a, big, female, giraffe, male, tall, the}\}$, $\start=\NT{Nominal}$ and define the following production rules:
\begin{figure}[t]
    \center
    \Tree [.$\NT{Nominal}$ [.$\NT{Det}$ a ] [.$\NT{NP}$ [.$\NT{N}$ giraffe ] ] ]
    \hspace{15mm}
    \Tree [.$\NT{Nominal}$ [.$\NT{Det}$ the ] [.$\NT{NP}$ [.$\NT{Adj}$ big ] [.$\NT{NP}$ [.$\NT{N}$ male ] ] ] ]
    \hspace{15mm}
    \Tree [.$\NT{Nominal}$ [.$\NT{Det}$ a ] [.$\NT{NP}$ [.$\NT{Adj}$ tall ] [.$\NT{NP}$ [.$\NT{Adj}$ female ] [.$\NT{NP}$ [.$\NT{N}$ giraffe ] ] ] ] ] 
    \caption{Derivation trees for natural language nominal phrases.}
    \label{fig:ex_np_grammar}
\end{figure}
\begin{align*}
    & \production{\NT{Nominal}}{\NT{Det}\; \NT{NP}}\\
    & \production{\NT{NP}}{\NT{N} \mid \NT{Adj}\; \NT{NP}} \\
    & \production{\NT{Det}}{\text{a}\mid \text{the}} \\
    & \production{\NT{N}}{\text{female} \mid \text{giraffe} \mid \text{male}}\\
    & \production{\NT{Adj}}{\text{big} \mid \text{female} \mid \text{male} \mid \text{tall}}
\end{align*}

See \cref{fig:ex_np_grammar} for a few examples of derivation trees in this grammar. Consider the tree with yield $\str=\text{a giraffe}$. Its left-most derivation is $$(\NT{Nominal},\NT{Det}\; \NT{NP}, \text{a}\; \NT{NP}, \text{a}\; \NT{N}, \text{a giraffe}).$$ Its right-most derivation is $$(\NT{Nominal}, \NT{Det}\; \NT{NP}, \NT{Det}\; \NT{N}, \NT{Det}\; \text{giraffe}, \text{a giraffe}).$$ Hence, we note that the mapping from derivations to derivation trees follows a many-to-one relationship---it is non-injective. 

%We provide this grammar in the notebook for Lecture 7. 

%\Ryan{I would put articles in the trees to make them more realistic.}
\label{ex:nominal_grammar}
\end{example}

\begin{example}[Recognizing $a^n b^n$]
The language $\lang=\{a^n b^n\mid n\in\mathbb N\}$ is not regular.
However, we can show that it is context-free with a simple grammar $\grammar$ with $\nonterm = \{X\}$, $\alphabet = \{a, b\}$, $\start=X$, $\rules=\{\production{X}{aXb}, \production{X}{\varepsilon}\}$.
\begin{proof}
We will show that $\lang=\lang(\grammar)$ in two steps: (i) showing that $\lang\subseteq\lang(\grammar)$ and (ii) showing that $\lang(\grammar)\subseteq\lang$. Define $\str_n=a^nb^n$.

(i) We first need to show that each $\str\in\lang$ can be generated by $\grammar$, which we will do by induction.
\paragraph{Base case ($n=0$)}
We have that $\str_0=\varepsilon$, which is generated by $\tree=(\production{X}{\varepsilon})$.
\paragraph{Inductive step ($n>1$)}
We have that $\str_n$ is generated by $$\tree=\underbrace{(\production{X}{aXb})\cdots(\production{X}{aXb})}_{n\text{ times}}(\production{X}{\varepsilon}).$$ 
It is then easy to see that $\str_{n+1}$ is generated by the derivation we get by replacing the last rule $(\production{X}{\varepsilon})$ with $(\production{X}{aXb})(\production{X}{\varepsilon})$. See \cref{fig:trees_example1} for an illustration.

(ii) Next, we show that for each $\tree\in\derivation{\grammar}$, we have that $\str(\tree)\in\lang$.
\paragraph{Base case ($\tree=(\production{X}{\varepsilon})$)}
It is trivial to see that the derivation $\tree=(\production{X}{\varepsilon})$ yields $\str(\tree)=\varepsilon$.

\paragraph{Inductive step}
Now observe that $\rules$ only contains two production rules and one non-terminal. Starting with $X$, we can either apply $\production{X}{aXb}$ to get one new non-terminal $X$, or apply $\production{X}{\varepsilon}$ to terminate the process. Hence, if we fix the length of the sequence of production rules there is no ambiguity in which string that will be generated. Thus, by induction, we conclude that if we have a derivation tree given by $\underbrace{(\production{X}{aXb}),\dots,(\production{X}{aXb})}_{n\text{ times}},(\production{X}{\varepsilon})$ generating $a^nb^n$, the derivation tree given by $\underbrace{(\production{X}{aXb}),\dots,(\production{X}{aXb})}_{n+1 \text{ times}},(\production{X}{\varepsilon})$ will generate $a^{n+1}b^{n+1}$.

\end{proof}
\begin{figure}[t]
    \center
    \Tree [.X $\varepsilon$ ]
    \hspace{10mm}
    \Tree [.X a [.X $\varepsilon$ ] b ]
    \hspace{10mm}
    \Tree [.X a [.X a [.X $\varepsilon$ ] b ] b ]
    \hspace{10mm}
    \Tree [.X a [.X a [.X a [.X $\varepsilon$ ] b ] b ] b ]
    \caption{A sequence of derivation trees for the strings in $\{a^nb^n \mid n=0,1,2,3\}$.}
    \label{fig:trees_example1}
\end{figure}

\label{ex:ab_grammar}
\end{example}

% \subsection{Derivation Sets and Ambiguity}

% We may have multiple derivation trees for a given string. Imagine, for instance, that we add a non-terminal $\NTY$ and rules $\production{X}{\NTY}, \production{\NTY}{\varepsilon}$ to the grammar in \cref{ex:ab_grammar}. The empty string $\varepsilon$ may then be derived either by $(\production{X}{\varepsilon})$, or $(\production{X}{\NTY}),(\production{Y}{\varepsilon})$, corresponding to two separate derivation trees. The set of these two trees comprise what we will call the \defn{derivation set} of the empty string. We will denote a derivation set of a string $\str$, generated by the grammar $\grammar$, as $\derivation{\grammar}{\str}$.

% We say that a grammar is \defn{unambiguous} if for every string that can be generated by the grammar, there is only one associated derivation tree. We define it formally using our newly introduced notation.

% \begin{definition}[Unambiguity]
% A grammar $\grammar$ is \defn{unambiguous} if $\forall\str\in\lang(\grammar)$, $|\derivation{\grammar}{\str}|=1$. 
% \end{definition}

% The converse holds for \defn{ambiguous} grammars.

% \begin{definition}[Ambiguity]
% A grammar $\grammar$ is \defn{ambiguous}  if $\exists\str\in\lang(\grammar)$ such that $|\derivation{\grammar}{\str}|>1$.
% \end{definition}

% We may also define ambiguity in terms of derivations. A grammar $\grammar$ is ambiguous if there exists a string in $\lang(\grammar)$ that has more than one left-most (or right-most) derivation. Conversely, a grammar is unambiguous if every string in its language has a unique left-most and right-most derivation.

% As in the FSA case, a subset of unambiguous context-free grammars are \defn{deterministic context-free grammars}. They are context-free grammars that can be derived from deterministic pushdown automata, which we will see in later lectures. %\andreas{Note: Deterministic CFGs seem to be defined from PDA, so intentionally leaving it a bit vague.} 

% \begin{example}[Ambiguous and Unambiguous Grammars]
%     Consider the context-free (and regular) language $\lang=\kleene{\alphabet}$ with $\alphabet=\{x\}$. We can generate $\lang$ with ambiguous CFGs, e.g. with a grammar $\grammar$ having
%     \begin{equation*}
%         \nonterm=\{X\}, \start=\NTX \text{ and } \rules=\{\production{\NTX}{x \NTX}, \production{\NTX}{\NTX x}, \production{\NTX}{\varepsilon}\}.
%     \end{equation*}
%     Note that this grammar is ambiguous since the trees $\tree_1=(\production{\NTX}{x \NTX}),(\production{\NTX}{\varepsilon})$ and $\tree_2=(\production{\NTX}{\NTX x}),(\production{\NTX}{\varepsilon})$ both yield the string $x$. We can, however, generate the same language with a grammar that is unambiguous, by removing one of the non-terminal production rules:
%     \begin{equation*}
%         \nonterm=\{\NTX\}, \start=\NTX \text{ and }\rules=\{\production{\NTX}{x \NTX}, \production{\NTX}{\varepsilon}\}.
%     \end{equation*}
    
%     \label{ex:ambiguous_unambiguous_grammars}
% \end{example}

\subsection{Weighted Context-free Grammars}

As in the case of finite-state automata, we will augment the classic, unweighted context-free grammars with real-valued weights.
We associate each rule $X \rightarrow \valpha$ with a weight $\weight=\productionWeight(\production{X}{\valpha})\in\R$.

\begin{definition}[Weighted Context-free Grammar]
A \defn{(real-)weighted context-free grammar} is a 5-tuple $\wcfgtuple$ where $\nonterm$ is a non-empty set of non-terminal symbols, $\alphabet$ is an alphabet of terminal symbols with $\nonterm \cap \alphabet = \emptyset$, $\start\in\nonterm$ is the designated start non-terminal symbol and $\rules$ is the set of production rules, where each rule $p \in \rules$ is of the form $\wproduction{\NTX}{\valpha}{\weight}$ with $\NTX\in\nonterm$, $\valpha\in\kleene{(\nonterm\cup \alphabet)}$ and $\weight=\productionWeight(\production{\NTX}{\valpha})\in\R$.
\end{definition}

Weights assigned to productions by WFCGs can be arbitrary real numbers.
Analogous to stochastic WFSAs (cf. \cref{def:stochastic-wfsa}) describing locally normalized language models, we define probabilistic WCFGs, where the weights of applicable production rules to a non-terminal form a probability distribution.
\begin{definition}
    A weighted context-free grammar $\grammar = \wcfgtuple$ is \defn{probabilistic} if the the weights of the productions of every non-terminal sum to $1$, i.e., for all $\NTX \in \nonterm$, it holds that 
    \begin{equation}
        \sum_{\production{\NTX}{\valpha}} \productionWeight\left(\NTX\right) = 1
    \end{equation}
\end{definition}

As should be clear, the one extension compared to an ordinary CFG is that each production rule is associated with a weight. 
Again analogously to the WFSA case, we say that a string $\str$ is in the language of WCFG $\grammar$ if there exists a derivation tree $\tree$ in $\grammar$ containing only non-zero weights with yield $\treeyield{\tree}=\str$. 

% We can transform any WCFG $\grammar=\wcfgtuple$ to an unweighted CFG $\grammar'=(\nonterm, \alphabet, \start, \rules, \productionWeight')$ by setting
% \begin{itemize}
%     \item $\productionWeight'(p)=1, \quad \forall \arule\in\rules$ where $\productionWeight(\arule)\neq\zero$
%     \item $\productionWeight'(p)=0, \quad \forall \arule\in\rules$ where $\productionWeight(\arule)=\zero$
% \end{itemize}
% We will refer to this transformation as \defn{booleanization}.

Similarly to how we defined weights of \emph{paths} in WFSAs, we can define a weighting over derivation trees and strings in the context of WCFGs. 
We start by giving a definition for the weight of a tree as the multiplication over all rules belonging to the tree, i.e. as a \emph{multiplicatively decomposable} function over the weights of its participating rules.

\begin{definition}
The weight of a derivation tree $\tree\in\derivation{\grammar}$ defined by a WCFG $\grammar$ is
\begin{equation}
    \productionWeight(\tree)=\prod_{(\production{\NTX}{\valpha})\in \tree}\productionWeight(\production{\NTX}{\valpha})
\end{equation}
\end{definition}
As in the WFSA case, we can also define the \defn{string sum} under a WCFG.
Intuitively, for a given string $\str$, the string sum is the sum of the weights of all derivations in the WCFG with yield $\str$.
\begin{definition}
The \defn{string sum} $\normConstant_{\grammar}(\str)$ of a string $\str$ generated by a WCFG $\grammar$ is defined by
\begin{align}
    \normConstant_{\grammar}(\str) & = \sum_{\tree\in\derivation{\grammar}{\str}}\productionWeight(\tree)\\
    & = \sum_{\tree\in\derivation{\grammar}{\str}}\prod_{(\production{\NTX}{\valpha})\in\tree}\productionWeight(\production{\NTX}{\valpha})
\end{align}
\end{definition}

\subsubsection{Normalizing Weighted Context-free Grammars}
As in the pathsum for WFSA, a \defn{treesum} is the sum of the weights of all the trees in a WCFG. 
We first define the treesum for symbols (non-terminals and terminals).

\begin{definition}[Treesum]
The \defn{treesum} for a non-terminal $Y$ in a grammar $\grammar$ is defined by
\begin{align}
 \normConstant_{\grammar}(Y) &= \sum_{\tree\in\derivation{\grammar}{Y}} \productionWeight(\tree) \\
 &=\sum_{\tree\in\derivation{\grammar}{Y}}\prod_{(\production{X}{\valpha})\in\tree}\productionWeight(\production{X}{\valpha})
\end{align}
The treesum for a terminal $a\in\alphabet\cup\{\varepsilon\}$ is defined to be
\begin{equation}
    \normConstant_{\grammar}(a) \defeq \one.
\end{equation}
\label{def:treesum}
\end{definition}
The treesum for a grammar is then simply the treesum for its start symbol.
\begin{definition}[Treesum WCFG]
The \defn{treesum} of a WCFG $\grammar=\wcfgtuple$ is
\begin{align}
    \normConstant_{\grammar} & =\normConstant_{\grammar}(\start)\\
    &= \sum_{\tree\in\derivation{\grammar}{\start}} \productionWeight(\tree) \\
    &=\sum_{\tree\in\derivation{\grammar}{\start}}\prod_{(\production{X}{\valpha})\in\tree}\productionWeight(\production{X}{\valpha})
\end{align}

%\begin{align}
% Z(\grammar) &= \sum_{\tree\in\derivation{\grammar}} \productionWeight(\tree) \\
% &=\sum_{\tree\in\derivation{\grammar}}\prod_{(\production{X}{\valpha})\in\tree}\productionWeight(\production{X}{\valpha})
%\end{align}
\end{definition}

When the grammar $\grammar$ we refer to is clear from context, we will drop the subscript and write e.g. $Z(\start)$. 

Although we can in some cases compute the treesum of a WCFG in closed-form, as we will see in the example below, we generally require some efficient algorithm to be able to do so.

\begin{example}[Geometric Series as a Treesum]
\label{ex:geometric_treesum}
    Consider the WCFG $\grammar=\wcfgtuple$, given by $\nonterm=\{\NTX\}$, $\alphabet=\{x\}$, $\start=\NTX$, and the rules:
    \begin{align*}
        & \wproduction{\NTX}{x\, \NTX}{1/2} \\
        & \wproduction{\NTX}{\varepsilon}{1}
    \end{align*}
    The language generated by $\grammar$ is $\lang(\grammar)=\{x^n\mid n\geq 0\}$. Further note that this grammar is unambiguous -- each string $\str=x^m$, for some $m\geq 0$, is associated with the derivation tree given by $\underbrace{(\wproduction{\NTX}{x\, \NTX}{1/3}),\dots,(\wproduction{\NTX}{x\, \NTX}{1/3})}_{m\text{ times}},(\wproduction{\NTX}{\varepsilon}{1})$. Due to the multiplicative decomposition over the weights of the rules, the weight associated with each derivation tree $\tree$ will hence be
    \begin{equation*}
        \productionWeight(\tree)=\left(\frac{1}{3}\right)^m\times 1 = \left(\frac{1}{3}\right)^m
    \end{equation*}
    Accordingly, we can compute the treesum of $\grammar$ using the closed-form expression for geometric series:
    \begin{align*}
        \normConstant_\grammar&=\sum_{m=0}^\infty\left(\frac{1}{3}\right)^m \\
        & = \frac{1}{1-1/3} \\
        & = \frac{3}{2}
    \end{align*}
\end{example}

Similarly to how we used pathsums in WFSAs to normalize globally normalized finite-state language models, we can use treesums to normalize weighted context-free grammars.
While we developed a specific algorithm for computing the pathsums in WFSAs, the derivation of a general treesum algorithm in WCFGs is more involved and beyond the scope of this course.\footnote{Again, the interested reader should have a look at the Advanced Formal Language Theory course.}


\begin{theorem}\label{thm:globally-locally-normalized-wcfgs}
    Let $\grammar$ be a globally normalized weighted context-free grammar. 
    Then, $\grammar$ can be locally normalized.
\end{theorem}


\subsection{Tightness of Context-free Language Models}

Again, an advantage of globally normalized context-free language models is that they are always tight.

\subsection{Single-stack Pushdown Automata}


As mentioned in \cref{sec:ch7}, FSAs can only recognize regular languages. To recognized \defn{context-free languages}, we must resort to \defn{pushdown automata}(PDA), a more general and more expressive type of automation.
The only difference between PDA and FSA is that PDA has an additional \emph{stack} for storing \emph{arbitrarily long} strings from a special alphabet, which allow the PDA to effectively work with unbounded memory.
We first give a formal definition of PDA.

\begin{definition}[Pushdown Automaton]
A \defn{pushdown automaton} is a 7-tuple $\pdatuple$ where
\begin{itemize}
    \item $\states$ is a finite set of states;
    \item $\alphabet$ is a finite set of input symbols (input alphabet);
    \item $\stackalphabet$ is a finite set of stack symbols (stack alphabet);
    \item $\stackstates \subseteq \stackalphabet$ is the set of initial stack symbols;
    \item $\initial \subseteq \states$ is the set of initial states; 
    \item $\final \subseteq \states$ the set of final or accepting states;
    \item $\trans \subseteq \states \times \left( \alphabet \cup  \set{\varepsilon} \right) \times \stackalphabet \times \states \times \stackalphabet^*$ a finite multi-set of transitions.
\end{itemize}
\end{definition}

In PDA, the transition function $\trans$ takes as argument not only the finite state, but also the current state of the stack, i.e. $\trans(\stateq, a, X)$, where $\stateq \in \states$, $a \in \alphabet$ and $X \in \stackalphabet$. The output of $\trans$ is a finite set of pairs $\left(p, \stackseq{}\right)$, where $p$ is a new state and $\stackseq \in \stackalphabet^*$ is a string of stack symbols that replaces $X$ at the top of the stack. For example, if $\stackseq = \varepsilon$, then the stack is popped and nothing is pushed onto the stack; if $\stackseq = X$, then the stack is unchanged; if $\stackseq = YZ$, then $X$ is popped and $YZ$ is pushed onto the stack. $X = \varepsilon$ means that the stack is neither read nor popped. We denotes each transition as $\pdaEdgenoweight{\stateq}{a}{p}{X}{\stackseq}$.

\begin{definition}
A \defn{stack} $\stack \in \stackalphabet^*$ is a finite sequence of stack symbols.
\end{definition}

\begin{definition}
A \defn{machine configuration} is a pair $\langle \stack, \stateq \rangle$ of a stack $\stack$ and a state $\stateq$.
At every transition step, PDA applies the transition function $\trans$ and transforms from one machine configuration to another where this transformation has the effect of ``outputting'' $a$ if $a \neq \varepsilon$.
\end{definition}


Now we introduce Weighted Pushdown Automata (WPDA). Similar to WFSA, WPDA is also a generalization of the PDA where the transitions are weighted with weights from a productionWeight. 
% Transitions are also sometimes denoted as $\trans(q, a, W) = (q)$

\begin{definition}[Weighted Pushdown Automaton]
A \defn{weighted pushdown automaton} $\pushdown$ is a 9-tuple $\wpdatuple$ over a productionWeight $\productionWeight=\semiringtuple$ where
\begin{itemize}
    \item $\states$ is a finite set of states;
    \item $\alphabet$ is a finite set of input symbols (input alphabet);
    \item $\stackalphabet$ is a finite set of stack symbols (stack alphabet);
    \item $\stackstates \subseteq \stackalphabet$ is the set of initial stack symbols;
    \item $\initial \subseteq \states$ is the set of initial states; 
    \item $\final \subseteq \states$ the set of final or accepting states;
    \item $\trans: \states \times \left( \alphabet \cup  \set{\varepsilon} \right) \times \stackalphabet \times \states \times \stackalphabet^* \rightarrow \semiringset$ a weighting function over the finite multi-set of transitions.
    \item $\initf : \initial \rightarrow \semiringset$ a weighting function over the set of initial state $\initial$;
    \item $\finalf : \final \rightarrow \semiringset$ a weighting function over the set of final states $\final$.
\end{itemize}
\end{definition}

In a WPDA, each transition $\trans$ is associated with a weight $w \in \semiringset$, of the  form $\pdaEdge{w}{q}{a}{p}{X}{\stackseq}$. We reuse the notation $\trans$ to denote the weighted multi-set of all transitions in $\pushdown$.
The transitions in WPDA can also be written in the form of $\trans(q, a, X) = \{(p, \stackseq; w), \cdots \}$ 


\begin{definition}
A \defn{probabilistic pushdown automaton (PPDA)} is a WPDA in which all weights are in the interval $[0,1]$ and for each pair of a stack symbol $X$ and a state $q$ the $\oplus$-sum of the weights of all transitions of the form $\pdaEdge{w}{q}{a}{p}{X}{\stackseq}$ equals to 1, i.e.
\begin{equation}
\sum_{\left\{ (\pdaEdge{w'}{q'}{a'}{p'}{X'}{\stackseq}) \in \trans \mid q'=q, X'=X \right\}} w' = 1
\end{equation}
\end{definition}


In the context of finite-state machines, we considered with the notion of a path and quantities based on paths, e.g. the path sum.
A natural question to ask at this point is, then, what is a good generalization of a path to the case of PDAs?
Of course, any generalization of a path to a PDA must take into account the fact that PDAs have stacks in addition to states.
In this course, we will define that generalization as a \defn{run}.

\begin{definition}
A \defn{run} $\arun$ in a PDA is an element of $\kleene{\trans}$ with \emph{consecutive transitions}, meaning that it is of the form $$\left( \pdaEdge{w}{q_1}{a_1}{q_2}{X_1}{\stackseq_1}, \pdaEdge{w}{q_2}{a_2}{q_3}{X_2}{\stackseq_2} \ldots  \pdaEdge{w}{q_{n-1}}{a_{n-1}}{q_n}{X_{n-1}}{\stackseq_{n-1}} \right),$$ where:
$a_i \in \alphabet \cup \{\varepsilon\}$ is the input symbol, $X_i \in \stackalphabet \cup \{\varepsilon\}$ is the symbol on top of the stack (being popped) and $\stackseq_i \in \stackalphabet^*$ is the string of stack symbols being pushed onto the stack and $w_i \in \semiringset$ is the corresponding weight of the transition.
\end{definition}

As in the case of path, the \defn{length} of a run is also the number of transition in it, denoted as $|\arun|$, and the \defn{yield} of a run is the concatenation of the input symbols on the edges along the path, denoted as $\yield\left(\arun\right)$. 
Furthermore, $\runs$ has the same definition as in \cref{sec:path-def}, i.e. sets of runs. 
Similarly, we use $\langle \prevq{\arun}$, $\prevs{\arun} \rangle$ and $\langle \nextq{ \arun }, \nexts{\arun} \rangle $ to denote the origin and the destination configurations of a run, respectively. 

\begin{definition}
A run is an \defn{accepting} run if it starts in an initial configuration and ends in a final configuration and has an empty stack at the end of the run.\footnote{In literature, e.g. \cite{sipser1996computation}, it is traditional to say a PDA can halt on final state or empty stack and they are proven equivalent. Here, we assume both from convenience.}
\end{definition}

\begin{definition}[Language of a PDA]
If $\pushdown$ is a PDA, then $L(\pushdown)$ is the set of strings $\str$ such that there exist an accepting run that yields $\str$. 
%We say $\pushdown$ that accepts the language $L(\pushdown)$.
\end{definition}

\begin{example}
We now give an example of a pushdown automaton $\pushdown$ that accepts the language $\{a^n b^n \mid n \in \mathbb{N}\}$ in \cref{fig:pda_example}. $\left( \pdaEdgenoweight{1}{\varepsilon}{2}{\varepsilon}{z_0},  \pdaEdgenoweight{2}{b}{3}{a}{\varepsilon} \right) $ is a run in $\pushdown$; $\left(\pdaEdgenoweight{1}{\varepsilon}{2}{\varepsilon}{z_0},  \pdaEdgenoweight{2}{b}{3}{a}{\varepsilon}, \pdaEdgenoweight{3}{\varepsilon}{4}{z_0}{\varepsilon} \right) $ is an accepting run in $\pushdown$.
\end{example}
\begin{figure}[ht] 
\centering
    \begin{tikzpicture}[node distance = 10mm]
    \tiny
    \node[state, initial] (q1) [] { $1$ }; 
    \node[state] (q2) [right = of q1] {$2$  }; 
    \node[state] (q3) [right = of q2] { $3$ }; 
    \node[state, accepting] (q4) [right = of q3] { $4$ }; 
    \draw[-{Latex[length=3mm]}] (q1) edge[bend left, above] node{ $\varepsilon, \varepsilon \rightarrow z_0$ } (q2) 
    (q2) edge[loop below] node{ $a, \varepsilon \rightarrow a$ } (q2) 
    (q2) edge[bend left, above] node{ $b, a \rightarrow \varepsilon$ } (q3) 
    (q3) edge[bend left, above] node{ $\varepsilon, z_0  \rightarrow \varepsilon$ } (q4) ;
    \end{tikzpicture}
\caption{The PDA that accepts the language $\{a^n b^n \mid n \in \mathbb{N}\}$.}
\label{fig:pda_example}
\end{figure}



\subsubsection{Converting CFGs to PDAs} \label{sec:cfg2pda}
One of the most prominent use cases of pushdown automata---both in NLP and in compiler theory---is the construction of linear-time parsing.
Recall from our discussion of parsing in previous lectures, that the exact parsing, e.g. using the CKY algorithm, generally runtimes in $\bigo{N^3}$ time where $N$ is the length of the sentence to be parsed.
PDAs allow us to construct linear-time parsing algorithm.
In the special case when we can represent a context-free grammar as a deterministic\ryan{Make sure this definition is given above}, we can exactly parse in linear time.
However, even for inherently ambiguous languages, it still beneficial to consider linear-time parsers.\ryan{Maybe add a bit on transition systems}
The general paradigm we will discuss for the creation of PDAs that serve as parsers is to convert a context-free grammar to a special PDA, in which there is only one state (except for the initial and accepting states) and there are only two types of transition: \shift{} and \reduce{}. We may include a special $\varepsilon$-transition $\pdaEdgenoweight{q}{\varepsilon}{q}{S}{\varepsilon}$ to ensure that the PDA halts with an empty stack. \footnote{An alternative choice is to define a unique ending stack symbol for PDA. This is the definition used in \rayuela{}.} We denote this PDA as \defn{Shift-Reduce WPDA}.
There are two general conversion strategies: the \defn{bottom-up} shift--reduce strategy and the \defn{top-down} shift--reduce strategy.
We discuss each in turn in the weighted case and give proofs of correctness.

\subsection{Multi-stack Pushdown Automata}


\subsubsection{Turing Completeness of Multi-stack Pushdown Automata}

Besides modeling more complex languages than finite-state language models from \cref{sec:finite-state}, (multi-stack) pushdown automata will also serve an important part in analyzing some modern language models that we introduce later.
Namely, we will show that Recurrent neural networks (cf. \cref{sec:rnns}) can \emph{simulate} any two-stack PDA.
This will be useful to reason about the computational expressiveness of recurrent neural networks because of a fundamental result in the theory of computation, namely, that two-stack PDAs are \emph{Turing complete}:
\begin{theorem}
    Let $\pushdown$ be a two-stack pushdown automaton. Then $\pushdown$ can simulate a Turing machine.
\end{theorem}
\begin{proof}
    See 
\end{proof}